\begin{Solbook}{1}
 \par \noindent  a) $B=\dfrac {3}{\sqrt {x} + 2}$ \hspace {2cm} b) $x=1$  \hfill \qedEX 
\end{Solbook}
\begin{Solbook}{2}
 \par \noindent  a) $x=5$ \hspace {2cm} b) $(2;1)$  \hfill \qedEX 
\end{Solbook}
\begin{Solbook}{3}
 \par \noindent Một học sinh lớp 9A trồng 3 cây, một học sinh lớp 9B trồng 4 cây \hfill \qedEX 
\end{Solbook}
\begin{Solbook}{4}
 \par \noindent  \immini { Vì hai dây nối tiếp xúc với khối cầu nên $AC$ và $BC$ lần lượt là hai tiếp tuyến của $(O)$ tại $A$ và $B$. Khi đó $$AC=BC=\sqrt {OC^2-OA^2}=\sqrt {18^2-10^2}=4\sqrt {14}\approx 15 \text { (m)}.$$ Như vậy, chiều dài của mỗi dây nối là khoảng $15$ m. }{ \begin {tikzpicture}[scale=1,line join=round,>=stealth] \def \r {1.5} \path (0,0) coordinate (O) (-155:\r ) coordinate (A) (-25:\r ) coordinate (B) (-90:\r ) coordinate (D) ($(B)!1!90:(O)$) coordinate (B') ($(A)!1!-90:(O)$) coordinate (A') (intersection of A--A' and B--B') coordinate (C); \draw (O) circle (\r ) (O)--(C)--(A)--(O)--(B)--(C); \foreach \x / \goc in {O/90,A/-155,B/-25,C/-90,D/-45} \draw [fill=white] (\x ) circle (1.5pt) ($(\x )+(\goc :3.5mm)$) node {$\x $}; \end {tikzpicture} }  \hfill \qedEX 
\end{Solbook}
\begin{Solbook}{5}
 \par \noindent  \begin {center} \begin {tikzpicture}[scale=1, font=\footnotesize , >=stealth] \def \r {2} \path (0,0) coordinate (O) (180:\r ) coordinate (A) (0:\r ) coordinate (B) (110:\r ) coordinate (M) (-110:\r ) coordinate (N) ($(M)!1cm!-90:(O)$) coordinate (m') (intersection of A--B and M--m') coordinate (C) ($(M)!1/2!(N)$) coordinate (H) ($(A)!1cm!-90:(M)$) coordinate (a') (intersection of A--a' and M--C) coordinate (E) (intersection of A--a' and M--N) coordinate (F) ; \draw (O) circle (\r ) (C)--(M)--(O)--(N)--(C)--(B)--(N)--(M)--(B) (M)--(A) (E)--(F) ; \par \foreach \x /\g in{O/45,A/-135,B/0,C/180,M/90,N/-90,H/135,E/90,F/-150} \draw [fill=white] (\x ) circle (1.5pt) +(\g :.3) node{$\x $}; \end {tikzpicture} \end {center} \begin {enumerate} \item $CO$ là đường trung trực của $MN$ nên $CO$ vuông góc với $MN$ tại $H$ \par Chứng minh $HO\cdot HC=MH^2 $ suy ra $ HO\cdot HC=\dfrac {1}{4}MN^2$ \item $\widehat {MAB}=90^\circ $ (chắn nửa đường tròn) \par $\widehat {EMA}=\widehat {OMB}$ (cùng phụ $\widehat {AMO}$) \par Suy ra $\widehat {EMA}=\widehat {MBA}$ \par Mà $\widehat {MBA}=\widehat {FMA} \,(\stackrel \frown {AM}=\stackrel \frown {AN})$ \par Suy ra $\widehat {EMA}=\widehat {FMA} $ \par Do $\triangle MEF$ có $MA$ vừa là đường cao, vừa là phân giác nên $\triangle MEF$ cân tại $M$ \end {enumerate}  \hfill \qedEX 
\end{Solbook}
\begin{Solbook}{6}
 \par \noindent  Chứng minh $x^2+y^2\ge 2xy\, (\text {với} \, x, y \ge 0)$ \par Từ đó suy ra $x^2+2y^2+3 \ge 2xy+2y+2, y^2 + 2z^2 + 3 \ge 2yz + 2z + 2$, $z^2 + 2x^2 + 3 \ge 2zx + 2x + 2$\\ \par $\dfrac {1}{x^2 + 2y^2 + 3} + \dfrac {1}{y^2 + 2z^2 + 3} + \dfrac {1}{z^2 + 2x^2 + 3} \le \dfrac {1}{2} \left ( \dfrac {1}{xy + y + 1} + \dfrac {1}{yz + z + 1} + \dfrac {1}{zx + x + 1} \right ) $\\ \par $= \dfrac {1}{2} \left ( \dfrac {1}{xy + y + 1} + \dfrac {1}{\dfrac {1}{x} + \dfrac {1}{xy} + 1} + \dfrac {1}{\dfrac {1}{y} + x + 1} \right ) = 0,5 \text { (do } xyz = 1 \text {)}$ \\ \par Dấu “=”, xảy ra khi $x = y = z = 1$  \hfill \qedEX 
\end{Solbook}
